\section{Background: The Design Space of Compiler IRs}
\label{sec:background}

Intermediate representations are the interface through which compiler components —
frontends, transformations, verifiers, and backends — cooperate. We surveyed the IRs
of production and research compilers across language, tensor, hardware, verification,
and database domains. Two structural observations motivate this paper.

\paragraph{Chains, not single IRs.} Every mature compiler we surveyed is a
\emph{chain} of IRs, each tuned to a stage: GCC lowers GENERIC trees to three-address
GIMPLE and then to register-transfer lists~\cite{merrill2003gimple}; GHC goes from
typed Core through STG to Cmm~\cite{peytonjones1992stg}; Rust passes through HIR, MIR,
and LLVM IR~\cite{mir}; Swift stages SIL before LLVM~\cite{swiftsil}; JavaScript
engines carry bytecode plus two or more optimizing graph IRs. A framework for
agent-driven compilation should therefore treat IR \emph{forms} as schemas and support
chains natively, rather than betting on one representation.

\paragraph{Forms are affordances, not expressiveness.} The classical forms are
largely inter-translatable: CPS, ANF, and SSA are known to correspond
\cite{kelsey1995ssa, appel1998ssa}, and sea-of-nodes IRs make scheduling implicit
where CFG-based IRs make it explicit~\cite{click1995sea, wurthinger2017graal}.
Expressiveness therefore does not discriminate among forms. What differs is what each
form makes cheap or awkward — identity stability, definition mechanics, verification
strength, and provenance (Table~\ref{tab:forms}) — precisely the properties that
matter when the ``user'' of the compiler is an automated agent rather than a C++ pass
author.

\begin{table}[t]
  \centering
  \small
  \begin{tabularx}{\textwidth}{@{}>{\raggedright\arraybackslash}p{2.5cm} >{\raggedright\arraybackslash}p{3.3cm} >{\raggedright\arraybackslash}p{2.5cm} >{\raggedright\arraybackslash}X@{}}
  \toprule
  \textbf{Form} & \textbf{Exemplars} & \textbf{Identity} & \textbf{Verification} \\
  \midrule
  Linear 3-address & GIMPLE~\cite{merrill2003gimple} & implicit indices & pass-local checks \\
  CFG + SSA & LLVM~\cite{lattner2004llvm}, MIR~\cite{mir}, SIL~\cite{swiftsil} & renumbered on print & verifier functions \\
  Sea of nodes & C2, Graal~\cite{click1995sea, wurthinger2017graal} & in-memory pointers & structural invariants \\
  Functional, let-bound & GHC Core/STG~\cite{peytonjones1992stg}, Core Erlang~\cite{coreerlang} & binding names & type systems \\
  Name-keyed graph & ONNX~\cite{onnx}, TF GraphDef & \textbf{stable names} & schema/shape checks \\
  Loop + schedule & Halide~\cite{ragankelley2013halide}, TensorIR~\cite{zheng2023tensorir} & names & shape checks \\
  Structured bytecode & wasm~\cite{haas2017wasm}, SPIR-V~\cite{spirvspec} & indices & load-time validation \\
  E-graph & egg~\cite{willsey2021egg}, TASO~\cite{jia2019taso} & hash-consed classes & rule-based \\
  \bottomrule
  \end{tabularx}
  \caption{IR forms and their affordances. Expressiveness is largely equivalent across
  forms; identity stability, definition mechanics, verification, and provenance differ
  — and those differences are what agents feel.}
  \label{tab:forms}
\end{table}

\paragraph{Identity.} Optimizing IRs in every surveyed production compiler renumber
their values on output: LLVM and MLIR print dense \texttt{\%0}, \texttt{\%1}, \dots
indices, MIR and SIL numbering shifts across passes, and bytecode indices move with
their tables. Only \emph{exchange} formats designed for serialization keep names —
TensorFlow's GraphDef and ONNX reference values by persistent string
identifiers~\cite{onnx} — and those formats carry no optimization infrastructure.
Renumbering is costly for an agent: textual references do not survive a pass boundary,
diffs are noisy, and cached observations about particular nodes are invalidated by
unrelated edits.

\paragraph{Definitions.} Operator definitions are overwhelmingly build-time artifacts:
TableGen and ODS generate C++ classes for LLVM and MLIR, Swift uses \texttt{.def}
files, and JVM bytecode semantics are fixed by specification~\cite{jvmspec}. Two
systems deviate in opposite directions: ONNX defines operators as versioned data
schemas with type constraints~\cite{onnx} but has no transformation engine, while TVM
registers operators at compiler runtime~\cite{chen2018tvm} inside a schedule-driven
compiler. No surveyed system combines data-defined definitions with first-class
transformations and verification.

\paragraph{Verification.} Strong verification precedents exist — the JVM bytecode
verifier~\cite{jvmspec}, WebAssembly's load-time type validation~\cite{haas2017wasm},
eBPF's static safety verifier~\cite{ebpf}, SPIR-V validation~\cite{spirvspec}, SIL's
per-stage verification~\cite{swiftsil}, and nanopass's automatically generated
per-stage verifiers~\cite{keep2013nanopass} — but all are static: they run at load
time or build time, none per transformation, and none are transactional.

\paragraph{Provenance.} Compiler IRs largely omit provenance; LLVM's optimization
remarks attribute transformations to source lines only weakly. By contrast, database
optimizers maintain a memo of equivalent expressions and remember which rewrite rules
derived each expression~\cite{graefe1995cascades}; this is the closest existing
analogue of the delta-and-journal model we adopt. Identity, data-defined definitions,
and provenance are the axes on which Brrr positions itself.
